% Source: https://artofproblemsolving.com/community/c3h1425736_a_big_problem_single_variable
% Stable identity: aops-topic-1425736-q01
\textbf{Single variable.} Let $x$ be the number of members. The conditions say $x=4p+1=6q+3=7r+4$ for some nonnegative integers $p,q,r$, i.e.
\[
x\equiv1\pmod4,\qquad x\equiv3\pmod6,\qquad x\equiv4\pmod7,\qquad 1\le x\le100 .
\]

\textbf{Combine the first two.} $x\equiv3\pmod6$ means $x$ is odd and $3\mid x$. Together with $x\equiv1\pmod4$ this is $x\equiv1\pmod4$ and $x\equiv0\pmod3$, i.e.\ $x\equiv9\pmod{12}$. (Conversely, $x\equiv9\pmod{12}$ gives $x\equiv3\pmod6$.)

\textbf{Add the third.} The candidates $9,21,33,45,57,69,81,93$ leave remainders $2,0,5,3,1,6,4,2$ on division by $7$. Only $81$ leaves remainder $4$.

So the band has $\boxed{81}$ members. Check: $81=4\cdot20+1=6\cdot13+3=7\cdot11+4$.
